\scanpage{31}
\mspage{29}
\uline{Chapitre} 3\deleted{\unsure{2|4}} \uline{Remarques sur les groupes} p-\uline{divisibles}.

3.1. Bien sûr, en même temps que pour les VA, je dois en rabattre sur les groupes formels. Je veux simplement signaler que la construction de $G(A)$ dans 2.6. doit pouvoir se paraphraser sans difficulté pour un groupe p-di\inserted{vi}sible $\phi$ sur un préschéma $S$ à caractéristiques résiduelles égales au même $p$. En effet les homomorphismes de ce groupe\deleted{s} dans le groupe formel associé à $\underline{G}_a$ sont triviales (en particulier, $\phi$ n'a pas d'automorphismes infinitésimaux). D'autre part (sauf erreur) le faisceau en \deleted{groupes} \inserted{modules} des $H^2$ de $\phi$ à valeurs dans le groupe formel \uline{a}ssocié à $\underline{G}_a$ (au sens du complexe du groupe $\phi$, variante formelle), est un Module localement libre de rang $g^*$, où $g^*$ est la dimension du groupe dual de $\phi$, soit $\phi^*$; plus précisément, ce Module doit être canoniquement isomorphe à $\underline{\operatorname{Lie}}(\phi^*)$\inserted{$=\underline{t}_{\phi^*}$}. \inserted{(}Tout ceci est suggéré par les résultats de Tate et Lubin sur les modules de groupes formels; je dois avouer d'ailleurs que je n'ai pas essayé de tirer au clair ce qu'il faut entendre, dans ce contexte, par ``le groupe formel associé à $\underline{G}_a$'', et s'il y a lieu de prendre le groupe formel purement infinitésimal.\inserted{)} Il s'impose alors de paraphraser la construction de Serre, en regardant l'extension \deleted{canonique} universelle de $\phi$ par un groupe formel vectoriel, qui sera ici le groupe formel (covariant) de $\underline{t}_{\phi^*}$. Désignant par $G(\phi)$ cette extension, son algèbre de Lie $H(\phi)$ sera une extension
\[
(*)\qquad 0\longrightarrow\underline{t}_{\phi^*}^{\vee}\longrightarrow\underline{H}(\phi)\longrightarrow\underline{t}_{\phi}\longrightarrow0.
\]
Bien entendu, $G(\phi)$ et par suite $\underline{H}(\phi)$ seront fonctoriels en $\phi$, et compatibles avec changement de base. Alors qu'il est certainement vrai, dans un sens qu'il conviendrait de préciser, que \deleted{\unsure{$\underline{H}(\phi)$}} $G(\phi)$ varie ``moins que $\phi$'', quand on fait varier $\phi$ infinitésimalement dans une famille, il n'est pas vrai pour autant que $G(\phi)$ soit indépendant de telles variations (à l'exception, probablement, de celles du premier ordre), i.e. puisse être considéré comme provenant d'un cristal en groupes \deleted{p-divisibles} \inserted{$\pm$ formels}. Cela sera le cas\deleted{,} seulement, sans doute, quand $\phi$ sera extension d'un groupe \deleted{\unsure{formel}} ind-étale par un groupe p-divisible torique, \inserted{et $S$ réduit à un point (?).} Dans le cas général, il semble intéressant d'étudier les variations
\margin{30}
